\documentclass[11pt,a4paper]{article}
\usepackage[T1]{fontenc}
\usepackage[utf8]{inputenc}
\usepackage{lmodern,microtype}
\usepackage{amsmath,amssymb,amsthm,mathtools}
\usepackage{booktabs,longtable,array,enumitem}
\usepackage{xcolor,geometry,fancyhdr,hyperref}
\geometry{margin=22mm}
\hypersetup{
  colorlinks=true,
  linkcolor=blue!55!black,
  urlcolor=blue,
  citecolor=blue!55!black,
  pdftitle={FIN ST1002--ST1091: Total Nadsoliton Ontological Persistence and Annihilation No-Go},
  pdfauthor={Krzysztof Żuchowski}
}
\pagestyle{fancy}
\fancyhf{}
\fancyhead[L]{FIN ST1002--ST1091}
\fancyhead[R]{Release 10.95}
\fancyfoot[C]{\thepage}
\setlength{\headheight}{14pt}
\setlength{\LTpre}{6pt}
\setlength{\LTpost}{6pt}
\newtheorem{theorem}{Theorem}
\newtheorem{proposition}[theorem]{Proposition}
\newtheorem{corollary}[theorem]{Corollary}
\newtheorem{definition}[theorem]{Definition}
\newcommand{\Mtot}{\mathcal M_{\mathrm{tot}}}
\newcommand{\nothingstate}{\dagger}
\newcommand{\status}[1]{\textsf{[#1]}}

\begin{document}

\begin{titlepage}
\centering
\vspace*{16mm}
{\Huge\bfseries FIN ST1002--ST1091\par}
\vspace{5mm}
{\LARGE Total Nadsoliton Ontological Persistence\par}
\vspace{2mm}
{\LARGE and the Annihilation No-Go\par}
\vspace{13mm}
{\Large Krzysztof \.Zuchowski\par}
{\normalsize Independent Researcher --- Fractal Information Theory Project\par}
{\normalsize ORCID: 0009-0002-0909-3613\par}
\vspace{10mm}
{\large Six research rounds --- 90 programmes\par}
{\large Release 10.95 --- 30 August 2026\par}
\vfill
\begin{minipage}{0.86\textwidth}\small
\textbf{Ontological scope.} In this report, one nadsoliton means the entire
fundamental being, not one localized excitation among several objects.  The
internal order remains nadsoliton $\to$ light $\to$ matter $\to$ emergent
observer.  The report derives no new physical law by naming the whole state a
``soliton.''
\end{minipage}
\vfill
{\small English research report.  Mathematical status labels are explicit;
no laboratory or external-audit evidence is claimed.}
\end{titlepage}

\begin{abstract}
This report corrects the type of the annihilation question in FIN.  If one
nadsoliton is the whole fundamental being, its annihilation cannot be modelled
as the decay of a localized particle inside a pre-existing space.  Four types
must be separated: the complete state, an internal pattern, a physical vacuum,
and literal nonexistence.  A vacuum is a state; literal nonexistence is not.

The central result is a conditional ontological annihilation no-go.  Let
$\Mtot$ be a nonempty state space for the entire nadsoliton, let
$\nothingstate\notin\Mtot$ denote nonexistence, and let
$\Phi_t:\Mtot\to\Mtot$ be a globally defined invariant evolution.  Then no
internal trajectory reaches $\nothingstate$.  If $\Mtot$ is normalized, it
also cannot reach the zero vector.  This is a theorem of state-space typing and
complete invariance, not a new repulsive force.  It does not explain why there
is something rather than nothing.

The frozen strict operator supplies two genuine corollaries.  The unitary
group $e^{-itA}$ preserves Hilbert norm.  The heat semigroup $e^{-tA}$
preserves total mass while exponentially erasing internal contrast.  Thus the
whole normalized state can persist while patterns, records or coherence die.
However, FIN has not yet derived the canonical total-state manifold, the full
nonlinear autonomous law, its global completeness, or a theorem excluding all
killing/boundary extensions.  Post-hoc normalization can hide extinction, so
normalization is protective only when it is dynamically conserved.

Bootstrap existence is shown not to imply stability.  An external pump is
inadmissible as a fundamental explanation when the nadsoliton is the entire
being; only internal zero-sum circulation is ontologically compatible.  An
internal observer cannot create a post-event record of total annihilation,
because the event would remove observer, apparatus, clock and record.  This is
an operational no-go, not proof of eternal existence.  Six model classes
falsify overstrong versions of the claim.  The report closes with a minimal
axiom/removal audit and programmes ST1092--ST1121 for constructing or refuting
the missing kernel-to-total-state bridge.
\end{abstract}

\tableofcontents
\newpage

\section{Executive answer}

The shortest scientifically defensible answer to ``what blocks the
nadsoliton from annihilation?'' is:
\begin{quote}
\emph{Under the stated one-nadsoliton ontology, nonexistence is not an
internal target state.  A complete invariant evolution of a normalized total
state therefore cannot reach it.  In the present strict channels the concrete
mathematical blockers are conservation of Hilbert norm for $e^{-itA}$ and
conservation of total mass for $e^{-tA}$.  These protect the total normalized
state, not every internal pattern.}
\end{quote}

The answer has three different epistemic levels:
\begin{enumerate}[label=\textbf{E\arabic*.}]
\item \status{Proven, conditional}: closure of a complete state space excludes
      a target outside that space.
\item \status{Proven}: the frozen finite strict operator generates a norm-
      preserving unitary group and a mass-preserving heat semigroup.
\item \status{Open}: FIN does not yet derive that these normalized spaces and
      channels constitute the unique full ontology and full dynamics of the
      total nadsoliton.
\end{enumerate}

This result is important mainly because it removes a category error.  It is
not yet a physical explanation of existence and it is not evidence for an
immortal object.

\section{Correction of the object type}

The previous ST912--ST1001 cycle correctly distinguished global norm,
localization, coherence and topological identity, but it treated the word
``nadsoliton'' partly as if it denoted a localized nonlinear excitation.  The
user's clarification changes the primary type:
\[
  \boxed{\text{one nadsoliton}=\text{the complete fundamental being}.}
\]
Consequently, localized structures, light, matter, observers and records must
be internal sectors or patterns of that one state.  Their disappearance is not
the disappearance of the carrier of all states.

\begin{longtable}{@{}p{.22\textwidth}p{.35\textwidth}p{.34\textwidth}@{}}
\toprule
Object & Mathematical type & What its disappearance means \\
\midrule\endhead
Total nadsoliton & complete state in $\Mtot$ & loss of the entire carrier/state space \\
Internal pattern & localization, excitation, correlation or subsystem state & redistribution, mixing or local decay \\
Vacuum & a distinguished physical state on an algebra & absence of selected excitations, not absence of being \\
Literal nothing & $\nothingstate$, no carrier/algebra/state & not an ordinary element of the physical state space \\
\bottomrule
\end{longtable}

The twelve FIN positions or modes must therefore not be read as twelve
fundamental nadsolitons.  They are degrees of freedom in a finite mathematical
representation of aspects of the one proposed total object.

\section{Minimal mathematical core}

\begin{definition}[Total-state model]
A total-state model is a pair $(\Mtot,\Phi)$ in which $\Mtot$ is a nonempty
state space and $\Phi_t:\Mtot\to\Mtot$ is the admitted evolution for every
admitted time $t$.  The evolution is \emph{complete} when it is defined for
all such times without a boundary exit or blow-up.
\end{definition}

\begin{definition}[Ontological annihilation extension]
Introduce a cemetery/nonexistence symbol $\nothingstate$ and the enlarged set
\[
  \widehat{\Mtot}=\Mtot\sqcup\{\nothingstate\}.
\]
An ontological annihilation occurs only if a separately specified law
$\widehat\Phi_t$ maps some $x\in\Mtot$ to $\nothingstate$.  Merely writing the
symbol does not create such a transition.
\end{definition}

\begin{definition}[Annihilation boundary]
For a representation in a larger linear or metric carrier, define
$\partial_0\Mtot$ as the limit points representing loss of normalization or
loss of the carrier.  On the Hilbert unit sphere and probability simplex the
zero vector lies at positive distance:
\[
 \inf_{\|\psi\|=1}\|\psi\|=1,
 \qquad
 \inf_{p\ge0,\,\mathbf1^Tp=1}\|p\|_1=1.
\]
Projective Hilbert space
$\mathbb P(\mathcal H)=(\mathcal H\setminus\{0\})/\mathbb C^*$ has no zero ray.
\end{definition}

\begin{theorem}[Total-nadsoliton annihilation no-go]
Let $\Mtot$ be the state space of the complete nadsoliton, let
$\nothingstate\notin\Mtot$, and let $\Phi_t:\Mtot\to\Mtot$ be globally
defined for all admitted times.  Then no internal trajectory reaches
$\nothingstate$.  If in addition there is a conserved functional
$Q:\Mtot\to(0,\infty)$ whose extension satisfies $Q(0)=0$, no nonzero orbit
reaches the zero state either.
\end{theorem}

\begin{proof}
By invariance, $\Phi_t(x)\in\Mtot$ for every $x\in\Mtot$ and every admitted
$t$.  Since $\nothingstate\notin\Mtot$, equality
$\Phi_t(x)=\nothingstate$ is impossible.  In the strengthened statement,
$Q(\Phi_t(x))=Q(x)>0$, whereas $Q(0)=0$, hence $\Phi_t(x)\ne0$.
\end{proof}

The proof is deliberately elementary.  Its strength is exact typing; its
weakness is that the decisive premises have not yet been derived from FIN's
kernel.  If ``total state'' is defined to be closed under evolution, the
result is close to tautological.  The scientific problem is therefore shifted
to constructing $\Mtot$ and proving global invariance/completeness without
assuming the desired conclusion.

\section{Round I: ST1002--ST1016 --- well-posedness}

The first round introduced the four-type ledger and tested the frozen strict
matrix.  Numerically, for the independently reconstructed $12\times12$
strict operator,
\[
 \|A-A^T\|_\infty=0,
 \qquad \|A\mathbf1\|_\infty=2.22\times10^{-16},
\]
\[
 \lambda_0\simeq-5.11\times10^{-16},\quad
 \lambda_1=0.7541211542070798,\quad
 \lambda_{\max}=2.3421820411463.
\]
For a normalized test vector, $e^{-3iA}$ preserved norm to displayed machine
precision.  For a normalized probability vector, $e^{-3A}$ preserved mass to
$2.2\times10^{-16}$ while reducing contrast from $0.1533110$ to $0.0126443$.
Thus internal homogenization and total-state persistence coexist.

The round also showed that a killed semigroup
\[
 T_t=e^{-t(A+\gamma I)}
\]
is a consistent counterextension.  Its survival mass is $e^{-\gamma t}$.
Therefore the strict no-killing result is not automatic for every possible
extension of $A$.

\section{Round II: ST1017--ST1031 --- boundary protection}

The annihilation boundary makes the missing premises visible.  Three exact
facts were established:
\begin{enumerate}
\item normalized Hilbert and probability states remain a positive distance
      from zero;
\item a positive conserved charge forbids zero;
\item for an autonomous locally Lipschitz ODE with $f(0)=0$, a nonzero
      trajectory cannot first hit zero in finite time without contradicting
      uniqueness.
\end{enumerate}

Neither regularity nor finite-time no-hit prevents asymptotic extinction:
\[
 \dot x=-x,\qquad x(t)=x_0e^{-t}\longrightarrow0.
\]
Dropping local Lipschitz regularity permits finite-time collapse:
\[
 \dot x=-\sqrt{x},\quad x(0)=1,qquad
 x(t)=\left(1-\frac t2\right)^2_+,qquad T_{\rm hit}=2.
\]
A useful sufficient barrier schema is a coercive functional satisfying
\[
 V(x)\to+\infty\quad(x\to\partial_0\Mtot),
 \qquad V(\Phi_t x)\le V(x).
\]
No such full FIN functional is presently strict-sourced.

\section{Round III: ST1032--ST1046 --- bootstrap and internal circulation}

Self-reference is not persistence.  The map $F(x)=2x$ has a fixed point at
zero but it is unstable.  The identity map has every point fixed and selects
nothing.  A contraction such as
\[
 F(x)=\tfrac12x+\tfrac12
\]
has a unique attracting fixed point, but contraction is an additional theorem
to prove, not a consequence of the word ``bootstrap.''

If the nadsoliton is the whole being, a fundamental external pump would place
a reservoir below or outside the whole and contradict the stated ontology.
The compatible replacement is an internal transfer
\[
  \sum_a J_a=0,
\]
or a closed circulation.  A three-sector witness used
\[
 J=\begin{pmatrix}0&1&-1\\-1&0&1\\1&-1&0\end{pmatrix},
 \qquad J^T=-J,\quad J\mathbf1=0.
\]
Then $e^{tJ}$ conserves both quadratic norm and the total sector sum.  This is
a mathematical example, not a strict derivation of FIN's internal resource
law.

Global persistence is compatible with subsystem irreversibility.  A global
pure Bell state has zero von Neumann entropy while either reduced subsystem
has entropy $\log2$.  Pinching preserves trace but removes accessible
coherence.  These facts support the distinction between the one total state
and its internally observed patterns; they do not establish a quantum
ontology for FIN.

\section{Round IV: ST1047--ST1061 --- the internal observer}

An operational completion must type observer $O$, apparatus $M$, environment
$E$ and record $R$ as internal patterns of the total state $\Omega$.  A record
is a correlation, so record erasure is weaker than annihilation of $\Omega$.

\begin{proposition}[No internal post-annihilation record]
If total annihilation removes the total state, it removes every internal
observer, apparatus, clock and record.  Hence no record created after that
event exists inside the theory.
\end{proposition}

This proposition does not prove that annihilation cannot happen.  It proves
that direct internal post-event verification is ill-typed.  Only model-
dependent precursors---loss of subsystem trace, gap closure, coherence decay
or approach to a boundary---can be recorded before the hypothetical event.

A physical vacuum remains operationally different from nothing: it is a state
on an observable algebra and may have correlations and response functions.
Nothing supplies no algebra and no probability distribution.

\subsection{Dual dynamics}

The same positive semidefinite $A$ supports
\[
 U_t=e^{-itA},\qquad P_t=e^{-tA}.
\]
The unitary channel preserves coherent norm and is reversible.  The heat
channel preserves mass while suppressing every nonzero spectral mode.  This
is a precise duality of functional calculus, not a proof that physical time
and diffusion time are the same, nor a Wick-rotation theorem.  It yields the
key persistence lesson:
\[
 \boxed{\text{the total normalized carrier may persist while internal
 information contrast becomes inaccessible.}}
\]

An internal oscillation may be a relational clock candidate through phase
differences $e^{-it(\lambda_j-\lambda_k)}$, but seconds require a dimensional
calibration, populated modes and a readout instrument.  Frequency alone does
not derive physical time.

\section{Round V: ST1062--ST1076 --- falsification}

Six model classes were compared.  They prevent any unconditional promotion of
the no-go.

\begin{longtable}{@{}p{.23\textwidth}p{.30\textwidth}p{.37\textwidth}@{}}
\toprule
Model class & Diagnostic & Verdict \\
\midrule\endhead
Closed normalized flow & constant norm/trace, complete orbit & total zero excluded conditionally \\
Sub-Markov killing & $\operatorname{Tr}\rho_t<1$ before renormalization & annihilation probability allowed \\
Stochastic cemetery jump & survival $S(t)=e^{-\gamma t}$ & requires hazard and $\nothingstate$ \\
Finite-time collapse & boundary hit with failed regularity/invariance & possible after a premise is removed \\
Topological sector & integer label while a gap remains open & protection fails at gap closure \\
External embedding & transfer from FIN sector to a larger complement & contradicts ``FIN is the whole'' if fundamental \\
\bottomrule
\end{longtable}

The sharpest falsification warning concerns renormalization.  If an
unnormalized surviving state has trace $s(t)<1$, replacing it by
\[
 \rho_t^{\rm cond}=\rho_t/\operatorname{Tr}\rho_t
\]
returns trace one and hides the lost survival weight.  Hence normalization
blocks annihilation only when it follows from the dynamical law; postselection
cannot be used as proof.

Projective state space has the same limitation: $[\psi]=[c\psi]$ for every
$c\ne0$, so it is blind to absolute survival amplitude.  A trace or survival
observable must accompany the ray description.

\section{Round VI: ST1077--ST1091 --- theorem, necessity and status}

\subsection{Strict corollaries}

For a finite self-adjoint $A$,
\[
 U_t^*U_t=I,
 \qquad \|U_t\psi\|=\|\psi\|.
\]
Thus nonzero Hilbert states do not become zero under the strict unitary group.
For $A\mathbf1=0$ and the admitted heat embedding,
\[
 \mathbf1^TP_tp=\mathbf1^Tp,
 \qquad
 \|P_t(p-u)\|_2\le e^{-\lambda_1t}\|p-u\|_2,
\]
where $u=\mathbf1/12$ and $\lambda_1=0.7541211542\ldots$.  Total mass
persists while internal contrast dies.

\subsection{Axiom ranking and removal audit}

\begin{longtable}{@{}p{.07\textwidth}p{.33\textwidth}p{.23\textwidth}p{.28\textwidth}@{}}
\toprule
Rank & Premise & Necessity & Removal result \\
\midrule\endhead
1 & nonempty typed total-state space $\Mtot$ & definitional & without it no object is modelled \\
2 & $\nothingstate\notin\Mtot$ & indispensable for internal no-go & add absorbing cemetery state \\
3 & globally defined invariant $\Phi_t:\Mtot\to\Mtot$ & indispensable dynamically & permit boundary exit or finite collapse \\
4 & conserved positive norm/mass/charge & independently testable strengthening & killed/sub-Markov evolution loses weight \\
5 & compactness, topology or coercive barrier & optional robustness & not minimal once closure is proved \\
\bottomrule
\end{longtable}

The minimal statement uses premises 1--3.  Premise 4 is valuable because it
can be checked directly and can prove zero exclusion even in a larger carrier.
Premise 5 may yield perturbative robustness but is not logically required.

\subsection{Strict versus conditional status}

\begin{longtable}{@{}p{.34\textwidth}p{.20\textwidth}p{.36\textwidth}@{}}
\toprule
Claim/object & Status & Boundary \\
\midrule\endhead
Unitary norm conservation & \status{Proven} & finite strict self-adjoint $A$ \\
Heat mass conservation & \status{Proven} & probability embedding and $A\mathbf1=0$ \\
Heat erasure of contrast & \status{Proven} & spectral gap of frozen $A$ \\
Total-state annihilation no-go & \status{Proven, conditional} & typed $\Mtot$ and complete invariant flow \\
Annihilation boundary $\partial_0\Mtot$ & \status{Constructed} & proposed formal device \\
Closed internal circulation & \status{Constructed} & toy witness, not strict-sourced \\
Canonical FIN total-state manifold & \status{Open} & no kernel-to-state theorem \\
Full nonlinear complete evolution & \status{Open} & no global well-posedness proof \\
No-killing theorem for all admissible laws & \status{Open} & only present channels checked \\
Internal operational model & \status{Open} & state/clock/preparation/instrument/environment/record incomplete \\
Physical eternal nadsoliton & \status{Not established} & no empirical or full dynamical proof \\
\bottomrule
\end{longtable}

\section{Relation to the kernels and FIN ontology}

This cycle uses the frozen strict generator only for the two spectral
corollaries.  It does not complete the legacy-to-strict bridge and does not
transfer any historical legacy physical role.  In particular, no claim about
electromagnetic coupling, gravity hierarchy, Planck scale, Standard Model or
cosmology follows.

The missing formal arrow is
\[
 (K_{\rm strict},A)
 \quad\not\Longrightarrow\quad
 (\Mtot,\Phi,\text{operational state})
\]
without additional construction.  An operator admits many state semantics and
many dynamical extensions.  The next research must derive or reject this
arrow, rather than naming one preferred semantics after the fact.

Fractal layers are most consistently treated, under the one-nadsoliton
ontology, as internal scales, refinements or coarse descriptions of the same
total state.  They are not additional fundamental beings below the
nadsoliton.  This permits a coarse observer to lose access to fine information
while the total state remains normalized.  It does not derive length, time,
energy, a light cone, light or matter.

\section{What has actually been learned}

\begin{enumerate}
\item The question ``why does the whole being not annihilate?'' is not the
      same as soliton collision stability.
\item Within a closed normalized model, literal nonexistence is absent by
      type.  This is exact but partly definitional.
\item Strict unitary norm and heat mass are real mathematical invariants.
\item Internal information patterns may disappear even when the total state
      persists; heat dynamics demonstrates this explicitly.
\item A bootstrap equation is not a stability theorem.
\item A fundamental external pump conflicts with the claim that the
      nadsoliton is the whole being; only internal balanced transfer is
      compatible.
\item Conditional normalization, projective quotienting and coarse observation
      can conceal survival loss.  The unnormalized trace must be audited.
\item Total annihilation cannot leave an internal post-event record.  This
      limits falsifiability but does not establish impossibility.
\end{enumerate}

\section{Recommended programmes}

The next round is ST1092--ST1106.  The percentages below rank probability of a
decisive \emph{local mathematical verdict}, not probability that FIN is a
physical theory.

\begin{longtable}{@{}p{.11\textwidth}p{.56\textwidth}p{.12\textwidth}p{.12\textwidth}@{}}
\toprule
Programme & Objective & Priority & Verdict probability \\
\midrule\endhead
ST1092 & classify canonical total-state candidates: sphere, rays, simplex, density states, algebraic states & 1 & 85\% \\
ST1093 & compare GNS, Dirichlet-form and spectral-state constructions from $A$ & 2 & 70\% \\
ST1094 & derive global existence/completeness criteria for every admitted nonlinear completion & 1 & 55\% \\
ST1095 & search for a strict coercive functional diverging at $\partial_0\Mtot$ & 2 & 40\% \\
ST1096 & audit whether normalization is a conserved observable or imposed gauge & 1 & 90\% \\
ST1097 & classify all positive/trace-preserving versus killed generator extensions compatible with $A$ & 1 & 90\% \\
ST1098 & construct a typed internal-sector decomposition without an external reservoir & 2 & 65\% \\
ST1099 & test whether $A$ sources a nontrivial zero-sum resource circulation & 3 & 35\% \\
ST1100 & build finite unitary dilations of effective subsystem heat/loss channels & 2 & 80\% \\
ST1101 & impose refinement compatibility on the total-state family & 2 & 55\% \\
ST1102 & include observer, apparatus, environment and record in one state map & 2 & 45\% \\
ST1103 & construct a relational clock and prove its domain/monotonicity limits & 3 & 55\% \\
ST1104 & classify compatibility, not identity, of unitary and heat dynamics & 1 & 85\% \\
ST1105 & automate counterexample search for boundary, killing and hidden postselection & 2 & 95\% \\
ST1106 & issue strict/conditional closure verdict for the kernel-to-total-state bridge & 1 & 95\% \\
\bottomrule
\end{longtable}

If ST1092--ST1106 produces a noncircular state map, ST1107--ST1121 should test
state-space uniqueness, exact global-existence intervals, boundary index,
coarse-layer consistency, internal-record theorems, vacuum sectors,
charge-versus-gauge ambiguity, refinement persistence, light/matter emergence
boundaries, falsification equivalence classes, dependency graphs, axiom
minimization and proof-assistant interfaces.  If it does not, the correct
outcome is a total no-go for deriving ontological persistence from the current
kernel alone.

\section{Final verdict}

\begin{quote}
\textbf{Deepest surviving interpretation.}  When one nadsoliton is defined as
the whole fundamental being, its literal annihilation is not an internal
dynamical process unless the theory first adds a larger state space, a
cemetery/nonexistence state, or a boundary/killing law.  Current strict FIN
contains no such law and its two principal channels conserve a global total.
Accordingly, present FIN blocks annihilation structurally---by normalized
state-space closure and conserved norm or mass---while allowing internal
patterns and accessible information to decay.  The missing theorem is not a
new force; it is a noncircular derivation of the total-state manifold and a
globally complete invariant full dynamics from the FIN kernel.
\end{quote}

This verdict survives the six countermodel classes.  It does not establish
that reality is FIN, that the nadsoliton exists, or that it persists physically
for infinite time.  It precisely states what would have to be proved before
such a conclusion could become mathematics rather than ontology by
definition.

\appendix

\section{Programme ledger}

\begin{longtable}{@{}p{.17\textwidth}p{.34\textwidth}p{.39\textwidth}@{}}
\toprule
Range & Principal object & Round verdict \\
\midrule\endhead
ST1002--ST1016 & total-being type ledger & vacuum/nothing separated; strict norm/mass checked \\
ST1017--ST1031 & annihilation boundary & finite/asymptotic countermodels isolate necessary premises \\
ST1032--ST1046 & bootstrap/internal circulation & fixed point not stability; external pump rejected fundamentally \\
ST1047--ST1061 & internal observer & no internal post-event record; dual dynamics separated \\
ST1062--ST1076 & falsification lattice & normalization, topology and totality stress-tested \\
ST1077--ST1091 & no-go theorem and synthesis & conditional theorem proved; strict kernel-to-state bridge open \\
\bottomrule
\end{longtable}

\section{Reproducibility}

The authoritative result sets are
\path{FIN_ST1002_ST1016_Results.json},
\path{FIN_ST1017_ST1031_Results.json},
\path{FIN_ST1032_ST1046_Results.json},
\path{FIN_ST1047_ST1061_Results.json},
\path{FIN_ST1062_ST1076_Results.json}, and
\path{FIN_ST1077_ST1091_Results.json}.  Every result references an individual
JSON packet by SHA-256.  The combined test
\path{test_fin_st1002_st1091.py} checks all 90 packet hashes and semantic
nonpromotion gates.

\section{Nonclaims}

No laboratory experiment, external audit, SI scale, physical light speed,
particle, photon, matter law, selector, QW-2191 discharge, legacy-to-strict
completion, legacy role transfer, gauge field, Standard Model, gravity,
$L_{\rm total}$ or Theory-of-Everything closure is established.  The report
does not prove that nonexistence is impossible outside the conditional model.
All future FIN PDFs remain English.

\begin{thebibliography}{9}
\bibitem{reed-simon}
M. Reed and B. Simon, \emph{Methods of Modern Mathematical Physics I:
Functional Analysis}, Academic Press.
\bibitem{engel-nagel}
K.-J. Engel and R. Nagel, \emph{One-Parameter Semigroups for Linear Evolution
Equations}, Springer.
\bibitem{perko}
L. Perko, \emph{Differential Equations and Dynamical Systems}, Springer.
\bibitem{nielsen-chuang}
M. A. Nielsen and I. L. Chuang, \emph{Quantum Computation and Quantum
Information}, Cambridge University Press.
\end{thebibliography}

\end{document}
